\section{训练基础设施}
本节介绍 Baichuan-M3 的训练基础设施，包括稳定的患者模拟环境、融合量表评估与事实感知评估的验证系统，以及渐进式的多阶段训练流水线。这些组件共同为长时程医学训练提供可靠的交互信号与可扩展的优化支撑。

\subsection{患者模拟器}
在我们之前的工作中~\cite{baichuan-m2}，引入了用于医患交互的患者模拟器。在部署过程中，我们发现当模拟主动型患者时模拟器会变得不稳定。具体而言，这类行为会打乱问诊流程，导致难以复现的模拟结果。为了在随机性带来的泛化收益与长时程训练所需的稳定性之间取得平衡，我们采用被动人格的患者模拟器，并实现两种互补的脚本驱动模式：被动交互模式与插入式打断模式。

\paragraph{被动交互模式（75\% 采样概率）：}
该模式仅提供患者画像、问诊量表与行为约束量表，不包含任何预定义对话历史。它模拟初诊的开场阶段。从空白交互状态开始，该模式评估医生智能体在高度不确定性下主动获取关键信息并形成初步诊断假设的能力。


\paragraph{插入式打断模式（25\% 采样概率）：}
该模式在基本患者信息上附加一段预定义对话片段，以模拟问诊中期状态。片段以患者发起的问题结束（通常与病情严重程度或治疗相关的焦虑问题），用于刻画问诊中的常见打断。然而，将该片段暴露给患者模拟器可能导致其模仿片段话风并偏离被动响应协议，从而引入不稳定性。因此我们采用非对称可见机制：片段仅对医生智能体可见，对患者模拟器隐藏。医生回答问题并继续问诊，模拟器被动接收并回应后续提问。

为降低模拟环境与真实在线问诊之间的分布偏差，我们在插入式打断模式中进一步引入细粒度概率配置。具体而言，半数样本为回合末提问，患者问题仅出现在片段末尾，以测试模型处理突发打断的能力；其余半数为回合中提问，在进行中的回合内插入问题以模拟多轮干扰。该混合配置确保模型在交互噪声（频繁打断、质疑与焦虑式追问）下仍具备稳健的诊断表现。


\subsection{验证系统}
为确保生成回复的可靠性与临床安全性，我们构建了综合验证系统，作为强化学习的主要奖励信号来源。不同于通用聊天模型通常只需流畅与有帮助，医疗智能体面临双重挑战：既要严格遵循诊疗规范，又要保证绝对的事实精确性。单一奖励模型往往难以解耦这些正交维度，可能鼓励“流畅的幻觉”而非严谨准确。为此，我们将评估过程解耦为两条并行流：\textbf{量表验证器}用细粒度标准评估结构质量与临床规范一致性，\textbf{事实验证器}则基于外部权威来源严格核验原子事实的生物学与医学有效性。这一混合方案确保模型同时优化专业合规与事实扎根。

\subsubsection{量表验证器}
M3 的验证栈沿用 M2 中提出的量表化验证范式~\cite{DBLP:journals/corr/abs-2508-12790,DBLP:journals/corr/abs-2507-17746,baichuan-m2}。我们不再将医疗回复质量视作单一偏好信号，而是把每次交互分解为一组可独立裁决的量表条款。基于 LLM 的评审器~\cite{DBLP:journals/corr/abs-2411-15594}逐项判定条款是否满足，并将结果聚合为标量奖励用于策略优化。

% Given a sample $x$ and a set of rubrics $\mathcal{R}=\{r_i\}_{i=1}^N$, each rubric $r_i$ is assigned a signed weight $w_i\in[-10,10]$, where $w_i>0$ denotes rewards and $w_i<0$ denotes penalties. An LLM judge outputs a binary decision $a_i\in\{0,1\}$ indicating whether $r_i$ is satisfied. The raw weighted score is $R_{\text{raw}}=\sum_{i=1}^N w_i a_i$. Let $R_{\max}$ and $R_{\min}$ denote the maximum and minimum attainable values of $R_{\text{raw}}$, which correspond to satisfying all reward rubrics or all penalty rubrics, i.e., $R_{\max}=\sum_{i:w_i>0} w_i$ and $R_{\min}=\sum_{i:w_i<0} w_i$. The task reward is obtained by min-max normalization:

给定样本 $x$ 与量表集合 $\mathcal{R}=\{r_i\}_{i=1}^N$，每条量表 $r_i$ 赋予带符号权重 $w_i\in[-10,10]$，其中 $w_i>0$ 表示奖励、$w_i<0$ 表示惩罚。LLM 评审器输出二值判定 $a_i\in\{0,1\}$ 表示 $r_i$ 是否被满足。任务奖励通过极差归一化得到：
\begin{equation}
R_{\text{task}}
=\frac{\sum_{i=1}^N w_i a_i-\sum_{i:w_i<0} w_i}{\sum_{i=1}^N |w_i|}
\end{equation}

该归一化将奖励尺度与量表数量及权重大小解耦到 $[0,1]$ 区间，从而使不同样本与量表配置下的奖励分布可比，同时可用整数权重以可解释的方式编码条款相对重要性。

为提高 RL 流水线效率，奖励评估与 rollout 异步调度：策略生成回复后，我们立即启动对应的量表评审任务，并将评审计算与其他互不冲突的 RL 操作（如计算 $\pi_\text{ref}/\pi_\text{old}$ 的对数概率）并行，从而加速训练。此外，我们为评审器采用前缀一致的提示词设计：系统约束、输出格式与对话上下文共享同一模板前缀，后缀替换为具体量表条款。该处理进行微批量化以最大化 KV-cache~\cite{DBLP:journals/tmlr/0002LTTXCHD0025} 复用，显著降低时间成本。

\subsubsection{事实验证器}
\label{sec:fact-verifier}
Baichuan-M2~\cite{baichuan-m2} 主要依赖量表化奖励来塑造医学推理。当在常见标准上趋于饱和时，模型往往通过添加更罕见的临床细节来追求边际收益，从而增加幻觉风险；因此我们同时优化完整性与可靠性。
为使“低幻觉率”的目标可量化、可优化，我们基于现有长文本事实性评估方法~\cite{DBLP:conf/emnlp/MinKLLYKIZH23,wei2024long}构建\textbf{事实感知验证流水线}。该流水线采用两阶段架构：首先将长文本回复拆解为细粒度、可独立验证的原子事实；其次使用检索增强的验证代理将每个事实与权威来源对照验证。该模块作为服务与量表验证器并行运行，尽可能降低时延。

\paragraph{原子事实抽取模型。}

该流水线的基础是将非结构化模型输出转换为离散、可验证的单元，称为原子事实。这些事实必须自洽，即使脱离原始上下文也应可被验证。为满足该要求，我们采用如下规则：

\begin{enumerate}
    \item \textbf{原子性与指代消解。} 将复杂复合句拆解为单事实单元，并解析代词以保证语义独立。
    \item \textbf{噪声与干扰项过滤。} 丢弃缺乏足够上下文的片段。尤其在选择题场景中，明确排除对错误选项（干扰项）的复述，避免被误判为幻觉。
    \item \textbf{保序去重。} 去除语义冗余断言，同时保持推理链的原始逻辑顺序。
\end{enumerate}

在评估阶段，我们最初使用 GPT-5~\cite{singh2025openai} 作为高质量抽取器；但其推理时延对在线 RL 过高。为此，我们从 GPT-5 蒸馏得到高效的 8B 抽取模型，使用涵盖多种医疗子任务的数据集训练，并作为最终抽取器；其细致的保真度评估见附录 \ref{sec:extraction_model_evaluation}。

\paragraph{检索增强验证器。}

抽取后，原子事实被输入检索增强验证代理。该代理对临床指南等权威医学来源进行迭代检索，并自主判断所收集证据是否足以支撑结论。每条事实最终被标注为 Supported、Refuted 或 Uncertain。相比仅依赖参数化知识或静态知识库的方法，这种动态检索机制可与最新临床证据保持一致，并更好地应对医学知识的持续演化。

\paragraph{两级缓存加速系统。}

将细粒度事实验证引入在线 RL 回路会带来显著的计算挑战。单次 RL 迭代可能产生成千上万条原子事实，对每条事实执行实时外部检索在成本与时延上都不可接受。

为加速训练，我们设计了两级事实缓存系统，基于如下关键观察：对于同一临床问题，模型多次采样回复在措辞上不同，但底层医学事实高度重合。

\begin{itemize}
    \item \textbf{一级缓存（精确匹配）。} 使用 Redis 缓存完全相同的事实字符串的验证结果，实现毫秒级检索与复用。
    \item \textbf{二级缓存（语义匹配）。} 将历史事实的嵌入存入向量数据库，使用 ANN 检索寻找语义等价的事实并复用验证结果。
\end{itemize}

随着缓存池增长，整体命中率从早期的不足 40\% 提升到约 80\%。这将外部检索请求减少约 85\%，使事实验证对训练的影响几乎可忽略。语义缓存不可避免会引入一定系统性偏差，例如剂量细微差异的事实可能被错误视为等价。我们通过第~\ref{sec:fact-aware} 节的信号去噪机制缓解该问题。

\subsection{多任务训练流水线}

为缓解多任务学习中的权衡问题~\cite{DBLP:journals/corr/abs-2007-10527,DBLP:conf/icassp/MoW22}并降低开发复杂度，Baichuan-M3 的训练流水线被拆分为三个渐进阶段：能力学习、分布融合与策略统一。通过将专家能力的获取与学生模型解耦，该设计实现了稳定的融合过程并显著提升整体性能。整体流水线如图~\ref{fig:overall} 所示。

\begin{figure}
    \centering
\includegraphics[width=1.0\linewidth]{image/Combine.pdf}
    \caption{三阶段流水线示意图。}
    \label{fig:overall}
\end{figure}

% 已经说明了全程的名词，不用再提一遍。
\paragraph{阶段 1：任务 RL。}

在初始阶段，我们的目标是构建多样且高质量的专家教师集合。基于共享初始化，我们为不同能力域部署独立的 RL 流水线。具体而言，为满足医疗应用的多样需求，我们显式训练临床问诊与健康咨询的专门专家，确保模型捕捉真实诊断对话与以患者为中心建议的细微差异。此外，我们还训练一个面向基础能力的通才专家，涵盖指令遵循与通用推理。

受 \cite{DBLP:journals/corr/abs-2510-13361} 启发，这一阶段的核心理念是“差异化而非统一”。允许不同模型在各自任务特定奖励下充分探索，我们获得一组具有强烈且不同归纳偏置的领域专家。该分治策略~\cite{DBLP:journals/corr/abs-2510-13361} 有效隔离任务间梯度干扰，避免多任务混合训练早期常见的优化冲突，从而为后续融合提供高置信行为指导。


\paragraph{阶段 2：离线策略蒸馏。}
第二阶段通过离线蒸馏将多个教师的能力“压缩”到单一学生模型中。我们冻结所有教师模型，并在各自领域内进行 rollout，构建离线轨迹数据集 $\mathcal{D}$。学生模型随后以离策略方式从该数据学习。

为适应离线数据的单样本现实并保证融合过程中的数值稳定性，我们用基于 Clip-Forward-KL 的受限蒸馏目标替代标准 KL 散度。对于每个样本 $(s, a)\in\mathcal{D}$ 及其对应的教师策略 $\pi_t$，损失函数定义为：

\begin{equation}
\mathcal{L}_{\text{clip-FKL}}(\theta) = \mathbb{E}_{(s, a) \sim \mathcal{D}} \left[ \mathbb{I} \left( \log \pi_\theta(a|s) < \log \pi_t(a|s) \right) \cdot \left( - \log \pi_\theta(a|s) \right) \right]
\end{equation}

其中 $\mathbb{I}(\cdot)$ 为指示函数，$\pi_\theta(\cdot)$ 与 $\pi_t(\cdot)$ 分别为学生与教师模型的逐 token 概率。该设计采用单边更新，仅在教师的经验支持上强制不劣性，而非过拟合整个条件分布。这避免了单样本 regime 下概率过度放大，并隐式保留数据支持之外的熵，有助于缓解模式塌陷，并为阶段 3 的寻模优化保留探索空间。

借助 Forward KL 的覆盖性特征，该阶段使学生模型广泛覆盖不同专家分布的高概率区域，从而稳定继承行为模式，消除直接多目标 RL 冷启动带来的不稳定性。

\paragraph{阶段 3：多教师在线策略蒸馏（MOPD）~\cite{mimo2025flash}。}
在第三阶段，学生模型重新进入在线交互环境，在混合领域分布上进行 rollout。此时模型同时受到真实任务奖励与多教师先验的约束。不同于第二阶段的分布模仿，这一阶段采用反向 KL 正则化~\cite{lu2025onpolicydistillation}。

利用 Reverse KL 的寻模特性，学生模型在多教师冲突建议下被驱动选择最优模式，而非被动平均。在真实奖励信号的引导下，学生从“模仿者”转变为“决策者”，在策略层面实现深度统一与去噪能力。

\paragraph{教师能力的演化。}
我们的框架并非静态流水线，而是支持循环式迭代优化。MOPD 得到的统一模型可作为阶段 1 的新初始化以进行领域特定增强，并随后进入新一轮蒸馏。该灵活性使模型能力得以在较低边际成本下持续提升。
